\ifdefined\iscombined
	\lecturetitle{Introduction}
	\setcounter{section}{1}
\else
\documentclass{article}
\usepackage{amsthm}
\usepackage{comment}
\usepackage{amsmath}
\usepackage{amsfonts}
\usepackage{sectsty}
\usepackage{hyperref}
\usepackage{fancyhdr}
\usepackage[top=1in,bottom=1in,left=1in,right=1in]{geometry}
\usepackage{float}
\usepackage{graphicx}
\usepackage{tabularx}
\usepackage{mathtools}

\title{Lecture 1 - Introduction}
\author{Matthew Farah, \href{mailto:farahm11@mcmaster.ca}{$\langle$farahm11@mcmaster.ca$\rangle$}, 400468251}
\date{\today}

\hypersetup{
        colorlinks=true,
        linkcolor=blue,
        filecolor=magenta,
        urlcolor=blue,
        citecolor=blue,
	anchorcolor=blue
}

\fancyhead{}
\fancyhead[L]{Matthew Farah, \href{mailto:farahm11@mcmaster.ca}{$\langle$farahm11@mcmaster.ca$\rangle$}, 400468251}
\fancyhead[R]{Lecture 1 - Introduction}

\newenvironment{directsubsubs}{
  \makeatletter
  \renewcommand\thesubsubsection{\thesection.\Roman{subsubsection}}
  \makeatother
}{
  \makeatletter
  \renewcommand\thesubsubsection{\thesubsection.\Roman{subsubsection}}
  \makeatother
}

\renewcommand{\thesubsection}{\thesection.\alph{subsection}}
\renewcommand{\thesubsubsection}{\thesection.\alph{subsection}.\Roman{subsubsection}}
\makeatletter

\sectionfont{\fontsize{12}{12}\bfseries}
\subsectionfont{\fontsize{10}{12}\normalfont}
\subsubsectionfont{\fontsize{10}{12}\normalfont}

\begin{document}

\maketitle
\pagestyle{fancy}

\fi

\begin{itemize}
	\item A \emph{sequential} system is a system composed of a single part which performs actions according to some defined order.
	\item A \emph{processing element} is defined as an environment in which a process is run, which is capable of allocating resources such as memory to the process running on it.
	\begin{itemize}
		\item By default, the word ``process'' refers to the execution of a sequential system.
	\end{itemize}
	\item \emph{Concurrent} and \emph{parallel} systems are systems composed of independently operating parts functioning together to perform some task.
	\begin{itemize}
		\item Concurrent and parallel systems are not the same thing, although they are often used interchangeably.
	\end{itemize}
	\item Concurrent systems consist of processes which perform actions with \underline{logical} simultaneity.
	\begin{itemize}
		\item Two concurrently running processes interleave their actions on the same processing element.
	\end{itemize}
	\item Parallel systems consist of processes which perform actions with \underline{physical} simultaneity.
	\begin{itemize}
		\item Two processes running in parallel require that each process run on a separate processing element.
	\end{itemize}
	\item Concurrent systems can be described as the composition of two or more sequential processes.
	\item There exist many ways to model concurrent systems:
	\begin{enumerate}
		\item \emph{Finite State Processes} (FSPs)
		\begin{itemize}
			\item Finite State Processes are an algebraic description of sequential and concurrent processes.
			\item Also called \emph{process algebras}.
		\end{itemize}
		\item \emph{Labelled Transition Systems} (LTSs)
		\begin{itemize}
			\item LTSs are extensions of finite state machines used to visually represent sequential and concurrent processes.
			\item \emph{Labelled Transition System Analyzers} (LTSAs) are tools used to visualize, animate, and analyze the behaviour of concurrent systems.
		\end{itemize}
		\item \emph{Petrinets}
		\begin{itemize}
			\item Petrinets are an alternate visualization of concurrent processes.
		\end{itemize}
	\end{enumerate}
\end{itemize}

\ifdefined\iscombined
\newpage
\else
  \end{document}
\fi
